\section*{Project contract}
Multiply an \(M\)-by-\(K\) matrix \(A\) by a \(K\)-by-\(N\) matrix \(B\).
Produce an \(M\)-by-\(N\) matrix \(C\), with
\[
  C_{i,j}=\sum_{k=0}^{K-1}A_{i,k}B_{k,j}.
\]
Each dimension is an integer from zero through sixteen. Input elements are signed
integers in \([-32768,32767]\), stored in four-byte words. Output elements are
exact signed integers stored in eight-byte doublewords. There is no fixed-point
scale, rounding, saturation or deliberate wraparound in this exercise.

On entry, \texttt{a0}, \texttt{a1} and \texttt{a2} hold the full addresses of
\(A\), \(B\) and \(C\). Register \texttt{a3} points to four consecutive words:
\(M,K,N,T\), in that order. The tile edge \(T\) is an integer from one through
eight. All matrices use packed row-major storage. Input pointers and the descriptor
are word aligned; the output pointer and incoming stack pointer are doubleword
aligned. Selected input, output and descriptor byte intervals are pairwise disjoint.
Do not assume the pointers equal the buffer labels or fit in thirty-two bits.

The active input lengths are \(4MK\) and \(4KN\) bytes; the active output length
is \(8MN\) bytes. Only the selected output bytes may change in caller-owned data.
Reads must remain in the selected slices and descriptor. Private stack storage
and submission-owned constants are allowed. Initialize an output byte before reading
it as a partial sum; preexisting output contents are arbitrary.

Return \(MN\) in \texttt{v0}. Preserve every bit of \texttt{s0} through
\texttt{s7}, \texttt{fp} and \texttt{sp}, and return through the supplied
\texttt{ra} continuation. Produce no console output. When \(M=0\) or \(N=0\),
perform no matrix reads or writes. When \(K=0\), write \(MN\) zero doublewords
without reading either input matrix. Descriptor and private-stack accesses remain
permitted in all three cases.

Keep the supplied caller, arena labels and allocation sizes for grading. The direct
checks may supply \texttt{function\_return}, a distinct continuation that rejoins the
caller after the function has returned. The actual caller returns to
\texttt{matmul\_done}. Use the supplied \texttt{ra} in either case.

Develop \texttt{matmul\_naive} with a row--column--reduction traversal, then
\texttt{matmul\_tiled} with clipped tiles. Both implement the same numerical and
ownership contract. Final free-writing grading accepts other correct schedules.
An instructor may separately assess the requested blocking strategy; numerical
agreement alone does not demonstrate that a program uses tiles.
